\documentclass[10pt]{article}
\providecommand{\FIGDIR}{fig}
\newcommand{\DIGESTDATE}{4 August 2026}
% ===================== PM Research Watch - shared preamble =====================
\usepackage[T1]{fontenc}
\usepackage[utf8]{inputenc}
\usepackage{amsmath,amssymb}
\usepackage{tgpagella}
\usepackage{tgheros}
\renewcommand{\sfdefault}{qhv}
\renewcommand{\ttdefault}{lmtt}
\usepackage{microtype}
\usepackage[english]{babel}

\usepackage[a4paper,top=17mm,bottom=17mm,left=17mm,right=17mm,
            headsep=6mm,footskip=10mm]{geometry}
\usepackage{graphicx}
\usepackage{xcolor}
\usepackage{colortbl}
\usepackage{tikz}
\usetikzlibrary{calc}
\usepackage[most]{tcolorbox}
\usepackage{titlesec}
\usepackage{fancyhdr}
\usepackage{booktabs}
\usepackage{array}
\usepackage{longtable}
\usepackage{enumitem}
\usepackage{multicol}
\usepackage{ragged2e}
\usepackage[hidelinks]{hyperref}

% ------------------------------------------------------------------ palette
\definecolor{Ink}    {HTML}{14262E}
\definecolor{Deep}   {HTML}{0F4C5C}
\definecolor{Teal}   {HTML}{1F7A8C}
\definecolor{Sky}    {HTML}{4FA3C4}
\definecolor{Sage}   {HTML}{6FA287}
\definecolor{Amber}  {HTML}{E3A008}
\definecolor{Clay}   {HTML}{B5643C}
\definecolor{Coral}  {HTML}{D1495B}
\definecolor{Violet} {HTML}{6A5B9E}
\definecolor{Slate}  {HTML}{4A6572}
\definecolor{Paper}  {HTML}{FBF9F5}
\definecolor{Mist}   {HTML}{EEEAE2}
\definecolor{Rule}   {HTML}{D9D3C9}

\pagecolor{Paper}
\color{Ink}

\hypersetup{colorlinks=true, linkcolor=Deep, urlcolor=Teal, citecolor=Teal,
            pdftitle={PM Research Watch}, pdfauthor={Automated daily digest}}

% ------------------------------------------------------------------ headings
\titleformat{\section}
  {\sffamily\bfseries\large\color{Deep}}{}{0pt}
  {}[{\vspace{1.2mm}\color{Amber}\titlerule[1.1pt]}]
\titlespacing*{\section}{0pt}{5.5mm}{2.8mm}

\titleformat{\subsection}
  {\sffamily\bfseries\normalsize\color{Teal}}{}{0pt}{}
\titlespacing*{\subsection}{0pt}{3.6mm}{1.4mm}

\setlist[itemize]{leftmargin=4.2mm, itemsep=1.1mm, topsep=1.3mm, parsep=0pt,
                  label=\textcolor{Teal}{\footnotesize$\blacktriangleright$}}

% ------------------------------------------------------------------ header/footer
\pagestyle{fancy}
\fancyhf{}
\renewcommand{\headrulewidth}{0pt}
\renewcommand{\footrulewidth}{0pt}
\fancyhead[L]{\sffamily\scriptsize\color{Slate}PM RESEARCH WATCH\;\textcolor{Rule}{\textbar}\;\DIGESTDATE}
\fancyhead[R]{\sffamily\scriptsize\color{Slate}Particulate matter monitoring \& health effects}
\fancyfoot[C]{\sffamily\scriptsize\color{Slate}\thepage}

% ------------------------------------------------------------------ components
% metric tile
\newcommand{\metric}[3]{%
  \begin{tcolorbox}[enhanced, colback=white, colframe=Rule, boxrule=0.5pt,
    arc=1.6mm, left=2.2mm, right=2.2mm, top=1.8mm, bottom=1.8mm,
    borderline west={1.7mm}{0pt}{#3}, nobeforeafter, valign=center,
    height=17mm, width=\linewidth]
    \RaggedRight
    {\sffamily\bfseries\fontsize{17}{17}\selectfont\color{#3}#1}\\[0.5mm]
    {\sffamily\fontsize{6.5}{7.4}\selectfont\color{Slate}#2}
  \end{tcolorbox}}

% pull-quote / key insight box
\newtcolorbox{keybox}[1][]{enhanced, colback=Mist, colframe=Mist,
  arc=1.6mm, left=3.2mm, right=3.2mm, top=2.4mm, bottom=2.4mm,
  borderline west={1.7mm}{0pt}{Amber}, #1}

% subtopic band
\newcommand{\band}[2]{%
  \vspace{3.2mm}\noindent
  \begin{tcolorbox}[enhanced, colback=#2, colframe=#2, arc=1.2mm,
    left=2.6mm, right=2.6mm, top=1.1mm, bottom=1.1mm, boxrule=0pt,
    width=\linewidth, nobeforeafter]
    {\sffamily\bfseries\small\color{white}#1}
  \end{tcolorbox}\vspace{1.6mm}}

% one paper entry: {short cite}{title}{journal}{doi}{body}
\newcommand{\paperentry}[5]{%
  \par\vspace{1.9mm}\noindent
  {\sffamily\bfseries\footnotesize\color{Ink}#1}\;%
  {\sffamily\fontsize{7}{8}\selectfont\color{Slate}#3}\\[0.4mm]
  {\footnotesize\itshape\color{Deep}#2}\\[0.6mm]
  {\footnotesize #4}\\[0.3mm]
  {\fontsize{6.6}{7.6}\selectfont\color{Slate}\texttt{doi:}\,\href{https://doi.org/#5}{\texttt{\detokenize{#5}}}}
  \par}

% caption style
\newcommand{\figcap}[1]{\par\vspace{1.0mm}
  {\fontsize{7.2}{8.6}\selectfont\color{Slate}\textsf{\textbf{Fig.}}\; #1\par}\vspace{1.5mm}}

\providecommand{\permil}{\textperthousand}

\setlength{\parindent}{0pt}
\setlength{\parskip}{0pt}
\renewcommand{\arraystretch}{1.18}


\begin{document}

% ==================================================================== MASTHEAD
\thispagestyle{empty}
\begin{tikzpicture}[remember picture, overlay]
  \fill[Deep] (current page.north west) rectangle
      ($(current page.north east)+(0,-67mm)$);
  \fill[Amber] ($(current page.north west)+(0,-67mm)$) rectangle
      ($(current page.north east)+(0,-68.4mm)$);
  \foreach \x/\y/\r in {22/12/0.9, 41/26/1.5, 63/9/0.7, 88/31/1.1, 112/17/1.9,
                        138/29/0.8, 159/13/1.3, 178/54/1.0, 196/21/1.6,
                        30/44/0.6, 74/57/1.2, 125/48/0.9, 168/6/0.8, 191/59/0.7,
                        52/18/1.0, 100/60/0.6, 148/20/1.1, 66/49/1.4, 205/45/0.8}
    \fill[Sky, opacity=0.30] ($(current page.north west)+(\x mm,-\y mm)$) circle (\r mm);
\end{tikzpicture}

\vspace*{4mm}
{\sffamily\fontsize{7.6}{9}\selectfont\color{Sky}%
 AUTOMATED DAILY LITERATURE DIGEST\;\textbullet\;PhD RESEARCH WATCH\par}
\vspace{2.2mm}
{\sffamily\bfseries\fontsize{27}{29}\selectfont\color{white}%
 Particulate Matter\par}
{\sffamily\fontsize{27}{29}\selectfont\color{Sky}%
 Monitoring \& Health Effects\par}
\vspace{5.0mm}
{\sffamily\fontsize{8.4}{10}\selectfont\color{white}%
 Issue 2026-08-04\;\;\textcolor{Amber}{\textbar}\;\;Entry window 4 August 2026\;\;%
 \textcolor{Amber}{\textbar}\;\;13 records in scope\par}
\vspace{18mm}

% -------------------------------------------------------------- metric strip
\noindent
\begin{minipage}[t]{0.190\linewidth}\metric{13}{RECORDS IN SCOPE\\(5 off-topic, 5 held)}{Deep}\end{minipage}\hfill
\begin{minipage}[t]{0.190\linewidth}\metric{6}{SUBTOPIC CLUSTERS\\REPRESENTED}{Teal}\end{minipage}\hfill
\begin{minipage}[t]{0.190\linewidth}\metric{6}{RECORDS FROM CHINA\\(46\% OF CORPUS)}{Sky}\end{minipage}\hfill
\begin{minipage}[t]{0.190\linewidth}\metric{2}{EFFECT ESTIMATES\\EXTRACTED}{Amber}\end{minipage}\hfill
\begin{minipage}[t]{0.190\linewidth}\metric{2}{REVIEWS /\\EVIDENCE SYNTHESES}{Clay}\end{minipage}

\vspace{6mm}

% -------------------------------------------------------------- signal of day
\section*{\sffamily\bfseries\color{Deep}Signal of the day}
\vspace{-1.5mm}{\color{Amber}\rule{\linewidth}{1.1pt}}\vspace{2.5mm}

\begin{keybox}
{\sffamily\bfseries\small\color{Deep}Aged Canadian wildfire smoke produced a
\emph{negative} \AA ngstr\"om exponent below 500\,nm --- the sign assumption that
underwrites almost every optical PM retrieval broke in the field.}\\[1.4mm]
\footnotesize
Masoom et al.\ (\textit{Atmos Chem Phys}) tracked short- and long-range plumes from the
2023 Canadian fire season across North America and Europe with combined remote sensing
and in situ measurement. In a roughly week-long episode late in the season they observed
\textbf{concave spectral curvature in AOD with the maximum above the 340\,nm minimum
wavelength}, driving the \AA ngstr\"om exponent negative below 500\,nm, together with
\textbf{accumulation-mode volume median diameters reaching $\sim$800\,nm} and
non-monotonic single-scattering albedo in submicron particles. Each of those three
features violates a default that low-cost optical sensing and AOD$\rightarrow$PM inversion
quietly rely on: that fine smoke is spectrally well-behaved and size-stable. The
practical reading for distributed sensing is narrow and concrete --- an optical instrument
calibrated on background urban aerosol has no valid response function during an
aged-smoke episode, and a fixed correction factor will fail in the \emph{direction} of
the correction, not merely in magnitude. This is the lesson the regional-calibration
literature reached from co-location statistics, arriving here from aerosol physics
instead.
\end{keybox}

\vspace{3mm}

% -------------------------------------------------------------- exec brief
\section*{\sffamily\bfseries\color{Deep}What changed today}
\vspace{-1.5mm}{\color{Amber}\rule{\linewidth}{1.1pt}}\vspace{2.5mm}

\footnotesize
\begin{itemize}\setlength{\itemsep}{1.1mm}
\item \textbf{The health axis is genuinely thin today.} PubMed returned ten health-axis
and two sensing-axis records for the 4 August entry date, five of them off-topic. The
sensing signal came almost entirely from Crossref-by-journal and Consensus, not PubMed
--- the structural gap noted since 28 July, unchanged. All six Crossref deposits carry a
4 August \texttt{created} date and every PubMed record an Entrez date of 4 August or
earlier, so nothing here belongs to a later issue.

\item \textbf{Network \emph{siting} is emerging as a strand distinct from calibration.}
Two independent groups treat sensor placement as an explicit optimisation rather than an
afterthought: Magdaong et al.\ combine street aspect ratio, traffic emissions and
irradiance into a suitability raster and reach 23--24\% population coverage with 40 nodes
in Quezon City; Colarusso et al.\ define spatial and temporal observability indices for
opportunistic vehicle fleets and co-design fixed-node placement against them. Calibration
research has matured far faster than siting research, and these are the first two papers
this watch has seen that make siting quantitative.

\item \textbf{Third consecutive day with no PM-specific ratio estimate carrying a
confidence interval.} The two estimates in the forest plot are the O$_3$ and NO$_2$
co-pollutant results from Russo et al.; that paper's PM$_{2.5}$ and SO$_2$ associations
were reported only as ``less consistent'', without numbers. This is a property of the
window, not a pipeline fault --- see the figure caption.

\item \textbf{Paternal preconception exposure is the more interesting reproductive
finding.} Russo et al.\ window exposure across the 74-day spermatogenesis cycle and
resolve it into mitosis, meiosis I--II, spermiogenesis and spermiation --- a sharper
exposure definition than the trimester blocks this literature normally uses. The signal
sits in mitosis and meiosis. Guo et al.'s review independently nominates spermatogenesis
stage I (70--90 days pre-collection) as the PM$_{2.5}$ vulnerability window and
traffic-derived metal constituents as the driver: convergent, though a review citing the
same primary literature is not independent evidence.

\item \textbf{Two mechanistic papers, both extra-pulmonary or barrier-focused, both with
rescue arms.} Shi et al.\ place RAGE upstream of
PERK/\allowbreak eIF2$\alpha$/\allowbreak ATF4/\allowbreak CHOP in
PM-aggravated COPD and rescue with FPS-ZM1 and 4-PBA; Qi et al.\ route PM$_{2.5}$ renal
fibrosis through alveolar-macrophage IL-6 suppressing ADCY3--CREB--PGC1$\alpha$ and rescue
by IL-6 blockade or ADCY3 overexpression. Both use intratracheal instillation, which
remains the dosimetric weak point of this whole literature.

\item \textbf{Five instrumentation-relevant records were found but not shipped.} Their
DOIs verify in Crossref but no abstract is retrievable from Crossref, Europe PMC or
Consensus, and the Elsevier and IEEE landing pages are client-rendered. They are logged
as held --- not rejected --- and deliberately kept out of \texttt{seen.json} so they can be
carried once indexed. Prajapati et al.\ on PM$_{2.5}$ hygroscopicity in the eastern
Indo-Gangetic Plain is the one that matters: RH dependence is directly load-bearing for
low-cost sensor correction.
\end{itemize}

\clearpage

% ==================================================================== ANALYTICS
\section*{\sffamily\bfseries\color{Deep}Corpus analytics}
\vspace{-1.5mm}{\color{Amber}\rule{\linewidth}{1.1pt}}\vspace{3mm}

\noindent\includegraphics[width=\linewidth]{\FIGDIR/f1_subtopics.png}
\figcap{Subtopic assignment of the 13 in-scope records, single-label by dominant
contribution. \textit{Sensing, forecasting \& instrumentation} leads with 4 (31\%); the
remaining five clusters carry one or two records each. Four of the ten standing
subtopics --- neuro/mental health, cardiovascular/metabolic, respiratory/allergic as a
primary label, and occupational/indoor --- are empty in this window.}

\vspace{2mm}
\noindent\includegraphics[width=\linewidth]{\FIGDIR/f2_design_endpoint.png}
\figcap{Left --- study architecture. Modelling and inventory work dominates (4/13), ahead
of measurement campaigns, reviews, cohort epidemiology and experimental toxicology at
two each. \textbf{All 13 records initially fell through to ``Other / mixed'' on this
run}; the design map has been extended and an unmapped-design warning added so this
failure is loud rather than silent in future (see run log). Right --- health endpoint by
organ system. Seven of thirteen records carry no human health endpoint at all --- four
monitoring or modelling, two emissions --- the highest non-health share this watch has
recorded.}

\vspace{2mm}
\noindent\includegraphics[width=\linewidth]{\FIGDIR/f3_metric_geography.png}
\figcap{Left --- particle metric. Only five records use PM$_{2.5}$ as the sole exposure and
one pairs it with PM$_{10}$; the remaining seven are unspeciated PM, particle
\emph{number}, INP concentration or column-integrated optical properties. That spread is
itself the finding --- mass concentration is not the measurand in over half of today's
corpus. Right --- geographic provenance. China contributes 6 of 13 (46\%). Southeast Asia
and the Eurasian Arctic each appear once; South Asia and sub-Saharan Africa are absent
again.}

\vspace{2mm}
\noindent\includegraphics[width=\linewidth]{\FIGDIR/f6_lifecourse.png}
\figcap{Life-course windows addressed; counts are non-exclusive, and only four of
thirteen records carry a human life-course window at all. The preconception entry is
paternal rather than maternal --- an unusual and welcome inversion of this literature's
default framing.}

\vspace{2mm}
\noindent\includegraphics[width=\linewidth]{\FIGDIR/f4_heatmap.png}
\figcap{Subtopic $\times$ study-architecture cross-tabulation. The persistent structural
feature is visible again: the \textit{Sensing} row contains no health-endpoint linkage,
and the mechanistic row is entirely experimental with no human biomarker bridge. With 13
records spread across six subtopics and six architectures the matrix is necessarily
sparse --- read it as a coverage map, not a distribution.}

\vspace{2mm}
\noindent\includegraphics[width=\linewidth]{\FIGDIR/f5_forest.png}
\figcap{\textbf{Both estimates plotted are co-pollutant results, not PM results.} No
record in this entry window reported a PM$_{2.5}$ or PM$_{10}$ ratio estimate with a
confidence interval --- the third consecutive issue for which that holds. Russo et al.\
report RR 0.71 (95\% CI 0.55--0.93) per 5\,ppb O$_3$ and RR 0.73 (0.50--1.05) per 5\,ppb
NO$_2$ for pregnancy per IUI cycle across the 74-day spermatogenesis window; their
PM$_{2.5}$ and SO$_2$ associations were characterised as less consistent and given no
numeric value. The plot is included so that the absence is recorded explicitly rather
than inferred from a missing figure, and so a downstream weekly or monthly rollup does
not read three sparse days as data loss.}

\clearpage

% ==================================================================== DIGEST
\section*{\sffamily\bfseries\color{Deep}Paper-level digest}
\vspace{-1.5mm}{\color{Amber}\rule{\linewidth}{1.1pt}}\vspace{1mm}

{\fontsize{7.4}{9}\selectfont\color{Slate}Ordered by cluster. Each entry gives the
methodological core and the finding that matters, not an abstract paraphrase.\par}

\band{Sensing, forecasting \& instrumentation\hfill 4 records}{Amber}

\paperentry{Masoom et al.\ 2026}
{Long-range transport of Canadian wildfire smoke to Europe in 2023: aerosol properties
and spectral features}
{Atmos Chem Phys}
{Combined column remote sensing and in situ measurement of short- and long-range plumes
across North America and Europe through the April--October 2023 fire season. The
late-season episode is the payload: \textbf{concave AOD spectral curvature with the
maximum above the 340\,nm minimum wavelength, giving negative \AA ngstr\"om exponents
below 500\,nm}; accumulation-mode volume median diameters near 800\,nm; and
non-monotonic single-scattering albedo, rising through the UV--visible then flattening or
falling at longer wavelengths, consistently across most stations. Limitation: a rare
$\sim$one-week event characterised observationally, with no closure between the retrieved
optics and a measured composition, so the cause of the curvature is inferred rather than
demonstrated. For distributed optical sensing the implication is that smoke-episode
response cannot be extrapolated from background calibration.}
{10.5194/acp-26-10801-2026}

\paperentry{Magdaong et al.\ 2026}
{Criteria and an algorithm for siting a low-cost IoT air-quality network near roads}
{ISPRS Ann Photogramm Remote Sens Spat Inf Sci}
{GIS multi-criteria analysis combining street aspect ratio, traffic emissions, global
horizontal irradiance and road proximity into a 0--6 suitability surface for Quezon City,
then a heuristic selecting nodes by joint suitability and population raster under a
minimum-spacing constraint. A 40-sensor configuration covers \textbf{1.27--1.35 million
residents, 23.0--24.4\% of the city}, stably across weighting schemes while keeping
spatial dispersion balanced. Limitation: coverage is defined geometrically against a
population raster with no validation against measured concentration fields, so
``covered'' means ``near a node'', not ``well estimated''. Still the clearest statement
this watch has seen that siting is an optimisation problem with a stated objective, and
the criteria transfer directly to any near-road deployment in a low- or middle-income
city.}
{10.5194/isprs-annals-x-5-w4-2025-315-2026}

\paperentry{Colarusso et al.\ 2026}
{Mobile urban sensing: spatio-temporal observability analysis and optimisation}
{Internet of Things}
{Fleet-aware framework for hybrid mobile--fixed sensing on an Edge--Fog--Cloud
architecture, evaluated on real data from three sensing vehicles covering 11.37\,km$^2$
in Benevento, Italy. Introduces spatial and temporal observability indices and a
value-of-fleet metric, then uses spatio-temporal interpolation plus what-if analysis to
quantify how fleet size, route redundancy, sampling density and data loss propagate into
reconstructed-map accuracy; a greedy co-design procedure places fixed nodes against the
residual. Limitation: a single city, a single small fleet, and interpolation-based ground
truth --- the observability indices are internally consistent but not externally
validated. The framing is the transferable part: opportunistic mobile coverage becomes a
measurable quantity rather than a qualitative claim.}
{10.1016/j.iot.2026.101907}

\paperentry{Li et al.\ 2026}
{Summertime ice-nucleating particle concentrations over the Eurasian Arctic seas}
{Atmos Chem Phys}
{Shipborne immersion-mode INP measurements from the Arctic Century Expedition across the
Barents, Kara and Laptev Seas and adjacent high-Arctic islands, August--September 2021.
$N_{\mathrm{INP}}$ was lower than at continental high-latitude sites, most clearly below
$-15\,^\circ$C, implying sparse mineral-dust INP; concentrations peaked in ice-free
water, marginal ice zones and near land, and were very low inside the ice pack, with the
maximum north of Novaya Zemlya under back-trajectories from the western Siberian coast.
The authors conclude sources are local to regional with little long-range transport
signature. Limitation: one summer and one cruise track, so the seasonal and interannual
envelope is unconstrained. Peripheral to PM exposure, but relevant to interpreting
particle counts in clean marine air, where the concentration floor is set by local
sources rather than by transport.}
{10.5194/acp-26-10835-2026}

\band{Exposure assessment \& modelling\hfill 2 records}{Amber}

\paperentry{Buhendwa et al.\ 2026}
{Modelling the coupled PM$_{2.5}$--ozone system in China: mechanisms, model biases,
perspectives}
{Environ Pollut}
{Review of the compound-pollution era in which PM$_{2.5}$ and O$_3$ share a precursor
pool (VOCs, NO$_x$, SO$_2$, NH$_3$) and single-pollutant control can worsen the
co-pollutant --- the ``pollution seesaw''. The useful content is the bias inventory:
chemical transport models \textbf{underestimate secondary organic aerosol by about 40\%},
with substantial spread across models, seasons and regions; daytime HONO sources and OH
reactivity are missing or understated; ammonia and episodic agricultural-burning
inventories are uncertain; nocturnal boundary-layer mixing is poorly represented. These
errors propagate into data assimilation and into machine-learning models trained on CTM
output. Limitation: a narrative synthesis with no quantitative pooling, so the 40\% figure
is a reported central tendency, not a meta-analytic estimate. Directly relevant if CTM
fields are used as priors or features in a sensor-fusion product.}
{10.1016/j.envpol.2026.128887}

\paperentry{Zhang et al.\ 2026}
{H$_2$O$_2$-driven sulfate formation at air--water interfaces: stepwise mechanism and
accelerated kinetics}
{Environ Sci Technol}
{Theoretical study of the dominant global route to particulate sulfate. At the air--water
interface the pre-reaction complex is favoured and the reaction proceeds stepwise ---
nucleophilic attack of HSO$_3^-$ on H$_2$O$_2$ to a HOOSO$_2^-$ intermediate, then
isomerisation to sulfate --- with a \textbf{rate-determining barrier of 4.1\,kcal/mol},
lower than in bulk, and with the rate-determining step itself shifting. Interfacial
electric field and partial solvation together account for \textbf{89\% of the barrier
reduction}. Limitation: entirely computational, with no experimental kinetic validation,
and the idealised interface is not a deliquesced particle surface. It matters because
rising H$_2$O$_2$ from wildfire activity makes this pathway more important, and because
sulfate is both a large PM$_{2.5}$ mass fraction and one of its more hygroscopic ---
therefore optically most humidity-sensitive --- components.}
{10.1021/acs.est.5c17902}

\band{Reproductive \& developmental\hfill 2 records}{Sky}

\paperentry{Russo et al.\ 2026}
{Critical windows of air-pollution exposure during spermatogenesis and IUI outcomes}
{Paediatr Perinat Epidemiol}
{Secondary analysis of the Folic Acid and Zinc Supplementation Trial: 592 male partners
in the Salt Lake City area contributing \textbf{1{,}223 intrauterine insemination cycles}
over nine months. Residential exposure came from Community Multiscale Air Quality models,
averaged over the 74-day spermatogenesis window before insemination and separately over
mitosis, meiosis I--II, spermiogenesis and spermiation, with generalised linear mixed
models adjusted for age, season, income and co-pollutants. Per 5\,ppb across the full
window, \textbf{O$_3$ RR 0.71 (95\% CI 0.55--0.93)} and \textbf{NO$_2$ RR 0.73
(0.50--1.05)} for hCG$+$ pregnancy, concentrated in mitosis and meiosis; live-birth
results were similar and no pattern emerged for pregnancy loss. \textbf{PM$_{2.5}$ and
SO$_2$ were less consistently associated and no numeric estimate was given.} Limitation:
an infertility-treatment population is not the general population, residential modelling
ignores occupational and commuting exposure, and co-pollutant models cannot cleanly
separate a traffic mixture. The stage-resolved exposure definition is the methodological
contribution worth borrowing.}
{10.1111/ppe.70185}

\paperentry{Guo et al.\ 2026}
{Environmental pollutants and testicular toxicity: basic evidence and translational
prospects}
{Reprod Toxicol}
{Narrative review of testicular-level evidence for microplastics, bisphenols and air
pollution, against a reported $>$50\% decline in global sperm concentration over five
decades. The PM-relevant claims: \textbf{spermatogenesis stage I, 70--90 days before
semen collection, is nominated as the PM$_{2.5}$ vulnerability window}, with metal
constituents of traffic-derived PM$_{2.5}$ as the primary toxicity driver. Separately,
microplastics are reported in human testicular tissue with a gut-dysbiosis $\rightarrow$
immune activation $\rightarrow$ blood--testis barrier pathway, and BCAT1-mediated
ferroptosis is offered as a bisphenol injury mechanism. Limitation: no systematic search
or quality appraisal, so the window and constituent claims inherit whatever selection
bias the cited primary literature carries. Read alongside Russo et al.\ --- convergent on
timing, but not independent evidence for it.}
{10.1016/j.reprotox.2026.109310}

\band{Mechanistic toxicology\hfill 2 records}{Violet}

\paperentry{Shi et al.\ 2026}
{PM$_{2.5}$ aggravates airway inflammation and barrier disruption in COPD mice via
RAGE-dependent PERK/eIF2$\alpha$/ATF4/CHOP}
{J Appl Toxicol}
{Cigarette-smoke-induced murine COPD model with subsequent PM$_{2.5}$ exposure. PM
worsened lung function, emphysema, inflammation, apoptosis and epithelial tight-junction
damage, and activated the PERK/eIF2$\alpha$/ATF4/CHOP endoplasmic-reticulum stress arm.
The causal claim rests on two rescue arms: the RAGE inhibitor FPS-ZM1 and the ER-stress
modulator 4-PBA each attenuated the phenotype, placing RAGE upstream of ER stress rather
than parallel to it. Limitation: instilled PM at a dose the abstract does not state, a
smoke-primed rather than naive background, and pharmacological inhibitors of imperfect
selectivity --- FPS-ZM1 in particular. The RAGE$\rightarrow$ER-stress ordering is the novel
part; the downstream inflammatory readouts are well-trodden.}
{10.1002/jat.70358}

\paperentry{Qi et al.\ 2026}
{Urban PM$_{2.5}$ worsens renal injury and fibrosis through a lung--kidney IL-6/ADCY3
axis}
{J Hazard Mater}
{Repeated intratracheal urban PM$_{2.5}$ at an environmentally relevant dose combined
with renal ischaemia--reperfusion injury in mice, plus tubular epithelial cell work.
PM exposure worsened renal dysfunction, tubular injury, fibrosis and lipid-droplet
accumulation; the proposed route is PM $\rightarrow$ systemic inflammation with alveolar
macrophages as the likely IL-6 source $\rightarrow$ suppression of ADCY3--CREB--PGC1$\alpha$
$\rightarrow$ impaired fatty-acid $\beta$-oxidation, increased glycolysis, mitochondrial
dysfunction and fibrosis. IL-6 blockade and ADCY3 overexpression each rescued, in vivo
and in vitro. Limitation: the renal-IRI co-insult makes this a second-hit model, so the
finding concerns PM aggravating existing injury rather than causing nephropathy; and
intratracheal instillation again overstates deposited dose rate. The extra-pulmonary axis
and the metabolic mechanism are both worth tracking.}
{10.1016/j.jhazmat.2026.143119}

\band{Burden, policy \& mitigation\hfill 2 records}{Coral}

\paperentry{Bu \& Sun 2026}
{PM emissions from China's iron and steel industry: patterns, drivers, implications}
{J Air Waste Manag Assoc}
{LMDI decomposition over 99 Chinese sites, separating stack from fugitive emissions into
five process factors: steel output, energy intensity, waste-gas generation coefficient,
stack concentration and fugitive emission factor. \textbf{National PM emissions fell
23.6\% from 2011 to 2016.} Stack concentration was the dominant inhibitor (to
$-$31{,}558\,t in 2015--16); the fugitive factor flipped from $+$5{,}037\,t (2010--11) to
$-$23{,}320\,t; the waste-gas coefficient pushed emissions up throughout
($+$19{,}516\,t in 2015--16); energy intensity contributed under 25\%. Regionally the
picture inverts --- South and Northwest China show annual increases, with Northwest
fugitive emissions reaching 232\% of the annual change. Limitation: the series ends in
2016, predating the ultra-low-emission retrofit programme entirely. The
stack-versus-fugitive split is what matters for siting monitors near industry: the two
have completely different spatial signatures.}
{10.1080/10962247.2026.2698603}

\paperentry{Zhang et al.\ 2026}
{Particle-number emissions from a hybrid electric vehicle across cold and warm starts}
{J Hazard Mater}
{A China-6 compliant HEV tested over real-driving-emissions and RTS95 cycles with fresh
and bench-aged aftertreatment. The particle-number size distribution was trimodal and
dominated by an accumulation mode peaking near \textbf{90\,nm}, with most emission from
urban and highway segments and \textbf{engine-start events, cold and warm alike, as the
dominant source}. Counter-intuitively, catalyst ageing \emph{improved} overall PN
removal: oxidation efficiency fell but gasoline particulate filter filtration rose as an
ash layer established. Optimising warm-start fuel enrichment cut PN further on the aged
system. Limitation: one vehicle, so this characterises a configuration, not a fleet.
Directly relevant to near-road sensing --- a 90\,nm mode driven by discrete start events
is largely invisible to mass-based instruments.}
{10.1016/j.jhazmat.2026.143156}

\clearpage

\band{Other clinical endpoints\hfill 1 record}{Slate}

\paperentry{Yang et al.\ 2026}
{Long-term exposure to multiple air pollutants and chronic pain in middle-aged and older
Chinese adults}
{J Pain}
{China Health and Retirement Longitudinal Study, 2011--2020: \textbf{10{,}378
participants free of chronic pain at baseline, 8{,}047 incident cases over a mean 7.39
years} --- an incidence approaching 78\%, itself the study's most important number and a
strong hint that the outcome definition is broad. Time-dependent Cox and accelerated
failure time models were combined with weighted quantile sum regression over the
five-pollutant mixture. PM$_{2.5}$, PM$_{10}$, SO$_2$ and CO were each associated with
higher risk and earlier onset, and \textbf{SO$_2$ was the predominant mixture
contributor}, not PM. Limitation: no interval estimates appear in the abstract,
self-reported pain invites differential misclassification, and in a highly correlated
mixture WQS weights identify the best-measured proxy as readily as the causal agent ---
read ``SO$_2$ drives it'' as a hypothesis.}
{10.1016/j.jpain.2026.106400}

\clearpage

% ==================================================================== TABLE
\section*{\sffamily\bfseries\color{Deep}Corpus register}
\vspace{-1.5mm}{\color{Amber}\rule{\linewidth}{1.1pt}}\vspace{2.5mm}

{\fontsize{6.7}{8.4}\selectfont
\setlength{\tabcolsep}{2.4pt}
\begin{longtable}{@{}>{\RaggedRight}p{25mm}>{\RaggedRight}p{41mm}>{\RaggedRight}p{27mm}%
>{\RaggedRight}p{23mm}>{\RaggedRight}p{18mm}>{\RaggedRight}p{15mm}@{}}
\toprule
\textsf{\textbf{First author}} & \textsf{\textbf{Focus}} & \textsf{\textbf{Journal}} &
\textsf{\textbf{Design}} & \textsf{\textbf{Metric}} & \textsf{\textbf{Region}}\\
\midrule
\endfirsthead
\toprule
\textsf{\textbf{First author}} & \textsf{\textbf{Focus}} & \textsf{\textbf{Journal}} &
\textsf{\textbf{Design}} & \textsf{\textbf{Metric}} & \textsf{\textbf{Region}}\\
\midrule
\endhead
\multicolumn{6}{@{}l@{}}{\cellcolor{Mist}\textsf{\textbf{\textcolor{Deep}{Neuro / mental health}}}}\\[0.6mm]
\href{https://doi.org/10.1001/jamanetworkopen.2026.25015}{Dales} & Wildfire-sourced PM$_{2.5}$ and ED visits for migraine and other primary headache & \textit{JAMA Netw Open} & Case-crossover & PM$_{2.5}$ (wildfire) & Canada\\
\href{https://doi.org/10.1093/ageing/afag215}{Glover} & Umbrella meta-review of urban environmental exposures and cognitive health & \textit{Age Ageing} & Umbrella review / meta-analysis & PM$_{2.5}$ & Global\\
\href{https://doi.org/10.1016/j.nsa.2026.107019}{Goossens} & In utero air pollution exposure and pre-/postnatal brain development & \textit{Neurosci Appl} & Narrative review & PM$_{2.5}$ / PM$_{10}$ / PM$_{\mathrm{coarse}}$ & Global\\
\href{https://doi.org/10.1016/j.isci.2026.116617}{Ma} & PM$_{2.5}$-induced neurotoxicity and depression: exposure risks and mechanisms & \textit{iScience} & Narrative review & PM$_{2.5}$ + components & Global\\
\href{https://doi.org/10.3760/cma.j.cn112150-20250804-00752}{Pan} & Maternal PM exposure and offspring neurodevelopment: effects and mechanisms & \textit{Zhonghua Yu Fang Yi Xue Za Zhi} & Narrative review & PM (bulk) & China\\
\addlinespace[0.9mm]
\multicolumn{6}{@{}l@{}}{\cellcolor{Mist}\textsf{\textbf{\textcolor{Teal}{Cardiovascular \& metabolic}}}}\\[0.6mm]
\href{https://doi.org/10.3760/cma.j.cn112338-20251120-00843}{Huang} & PM$_{2.5}$ and dyslipidaemia in rural Chinese children; dietary effect modification & \textit{Zhonghua Liu Xing Bing Xue Za Zhi} & Cross-sectional (multistage) & PM$_{2.5}$ & China\\
\href{https://doi.org/10.32394/pe/218966}{Kazmierska} & Air pollution and cardiovascular disease: review of current data & \textit{Przegl Epidemiol} & Narrative review & PM$_{2.5}$ / PM$_{10}$ & Global\\
\href{https://doi.org/10.1007/s00125-026-06814-2}{Xi} & Air pollution, built environment and temperature vs. prediabetes phenotypes & \textit{Diabetologia} & Prospective cohort & PM$_{2.5}$ & Germany\\
\href{https://doi.org/10.1161/JAHA.125.046917}{Yang} & Short-term PM$_{2.5}$ and hospitalisation/ED use in atrial fibrillation & \textit{J Am Heart Assoc} & Case-crossover & PM$_{2.5}$ (wildfire-influenced) & USA\\
\addlinespace[0.9mm]
\multicolumn{6}{@{}l@{}}{\cellcolor{Mist}\textsf{\textbf{\textcolor{Sky}{Reproductive \& developmental}}}}\\[0.6mm]
\href{https://doi.org/10.1016/j.jes.2025.10.043}{Chen} & Prenatal air pollution, fetal growth trajectory phenotypes and child neurodevelopment & \textit{J Environ Sci (China)} & Prospective cohort & PM$_{2.5}$ / PM$_{10}$ & China\\
\href{https://doi.org/10.1016/j.ecoenv.2026.120566}{Huang} & Joint prenatal bisphenol + air-pollutant mixture exposure and preschool IQ & \textit{Ecotoxicol Environ Saf} & Prospective cohort (mixtures) & PM$_{2.5}$ / PM$_{10}$ & China\\
\href{https://doi.org/10.1007/s40572-026-00552-8}{Peterson} & Ambient air pollution, PAHs and pubertal development: critical review & \textit{Curr Environ Health Rep} & Critical review & PM$_{2.5}$ components & Global\\
\href{https://doi.org/10.1016/j.envres.2026.125313}{Xie} & Long-term PM$_{2.5}$ and timing of pubertal development in boys and girls & \textit{Environ Res} & Prospective cohort & PM$_{2.5}$ & China\\
\addlinespace[0.9mm]
\multicolumn{6}{@{}l@{}}{\cellcolor{Mist}\textsf{\textbf{\textcolor{Sage}{Respiratory \& allergic}}}}\\[0.6mm]
\href{https://doi.org/10.1007/s00484-026-03289-z}{Bayar Muluk} & Climate change and otolaryngology: airway disease impacts and emission reduction & \textit{Int J Biometeorol} & Narrative review & PM$_{2.5}$ & Global\\
\href{https://doi.org/10.4103/mgr.MEDGASRES-D-25-00164}{Liao} & Short-term air pollution, respiratory symptoms and lung function in asthma & \textit{Med Gas Res} & Panel study & PM$_{2.5}$ / PM$_{10}$ & China\\
\href{https://doi.org/10.3390/epidemiologia7040097}{Sanchez de Prada} & Ambient PM and asthma-coded admissions in children 0-3 y, Spain & \textit{Epidemiologia} & Ecological panel & PM$_{2.5}$ / PM$_{10}$ & Spain\\
\href{https://doi.org/10.1016/j.jhazmat.2026.143076}{Yang} & Allergenic vs. toxic risk of airborne seed hairs compared with ambient PM$_{2.5}$ & \textit{J Hazard Mater} & Field measurement & PM$_{2.5}$ + bioaerosol & China\\
\addlinespace[0.9mm]
\multicolumn{6}{@{}l@{}}{\cellcolor{Mist}\textsf{\textbf{\textcolor{Violet}{Mechanistic toxicology}}}}\\[0.6mm]
\href{https://doi.org/10.1016/j.jes.2025.12.074}{Ding} & PM$_{2.5}$ promotes vascular calcification via miR-27b-5p/HES1 signalling & \textit{J Environ Sci (China)} & In vitro + organoid + murine & PM$_{2.5}$ & China\\
\href{https://doi.org/10.1007/s00360-026-01698-5}{Maraschi} & Osmoregulatory disruption and metal bioaccumulation from settleable atmospheric PM & \textit{J Comp Physiol B} & Controlled exposure (ecotox) & Settleable APM & Brazil\\
\href{https://doi.org/10.1016/j.jes.2025.11.018}{Wu} & Macrophage NLRP3 inflammasome in real-ambient PM-induced hepatic glucose dysregulation & \textit{J Environ Sci (China)} & Real-ambient murine + co-culture & Real-ambient PM & China\\
\addlinespace[0.9mm]
\multicolumn{6}{@{}l@{}}{\cellcolor{Mist}\textsf{\textbf{\textcolor{Clay}{Occupational \& indoor}}}}\\[0.6mm]
\href{https://doi.org/10.3389/fmedt.2026.1796828}{Chen} & Nested suction instrument to cut intraoperative surgical-smoke PM$_{2.5}$ & \textit{Front Med Technol} & Bench / device experiment & PM$_{2.5}$ (surgical smoke) & China\\
\href{https://doi.org/10.1093/annweh/wxag055}{Cole} & Determinants of respirable crystalline silica exposure in Australian tunnelling & \textit{Ann Work Expo Health} & Observational exposure study & Respirable dust / RCS & Australia\\
\href{https://doi.org/10.1016/j.jhazmat.2026.143081}{Ding} & Stage-resolved hair Zn-isotope dynamics as a firefighting combustion-exposure tracer & \textit{J Hazard Mater} & Repeated-measures biomarker & Combustion-derived PM & China\\
\href{https://doi.org/10.3390/jox16040119}{Ramsden} & Xenobiotic hazards in aircraft cabin air & \textit{J Xenobiot} & Review + risk economics & Ultrafine particles & Global\\
\addlinespace[0.9mm]
\multicolumn{6}{@{}l@{}}{\cellcolor{Mist}\textsf{\textbf{\textcolor{Amber}{Sensing, forecasting \& instrumentation}}}}\\[0.6mm]
\href{https://doi.org/10.1016/j.jenvman.2026.130515}{Shi} & Adaptive spatiotemporal graph transformer for multi-site multi-step PM$_{2.5}$ forecasting & \textit{J Environ Manage} & Deep-learning model & PM$_{2.5}$ & China / India\\
\href{https://doi.org/10.3390/bios16070380}{Tsuchiya} & Seasonal and indoor-outdoor characteristics of fluorescent aerosol particles in offices & \textit{Biosensors} & Real-time sensor campaign & Bioaerosol / coarse & Japan\\
\addlinespace[0.9mm]
\multicolumn{6}{@{}l@{}}{\cellcolor{Mist}\textsf{\textbf{\textcolor{Amber}{Exposure assessment \& modelling}}}}\\[0.6mm]
\href{https://doi.org/10.1029/2025GH001621}{Solorzano-Araque} & Multiscale exposure modelling: LOTOS-EUROS + mobility surveys + particle toxicity & \textit{GeoHealth} & CTM + mobility integration & PM$_{2.5}$ / PM$_{10}$ & Colombia\\
\addlinespace[0.9mm]
\multicolumn{6}{@{}l@{}}{\cellcolor{Mist}\textsf{\textbf{\textcolor{Coral}{Burden, policy \& mitigation}}}}\\[0.6mm]
\href{https://doi.org/10.1002/hsr2.72886}{Karamzad} & Ischaemic heart disease burden in MENA 1990-2021; PM attributable fraction 28.2\% & \textit{Health Sci Rep} & GBD secondary analysis & Ambient PM & MENA\\
\href{https://doi.org/10.1371/journal.pone.0347989}{Rees} & HEPA filtration to reduce respiratory infection in care homes & \textit{PLoS One} & RCT process evaluation & Filtered PM & UK\\
\href{https://doi.org/10.1016/j.jes.2026.01.009}{Shan} & Policy-driven emission reductions and post-pandemic rebound, high-resolution inventories & \textit{J Environ Sci (China)} & Emission inventory & PM emissions & China\\
\addlinespace[0.9mm]
\multicolumn{6}{@{}l@{}}{\cellcolor{Mist}\textsf{\textbf{\textcolor{Slate}{Other clinical endpoints}}}}\\[0.6mm]
\href{https://doi.org/10.3390/dj14070440}{Lupita} & PM$_{2.5}$, preventive-care capacity and caries/edentulism burden across the EU & \textit{Dent J} & Ecological panel & PM$_{2.5}$ & EU-27\\
\href{https://doi.org/10.1016/j.jdin.2026.06.008}{Park} & Long-term PM exposure and risk of androgenetic alopecia & \textit{JAAD Int} & Nationwide cohort & PM$_{2.5}$ / PM$_{10}$ & South Korea\\
\href{https://doi.org/10.3171/2026.3.PEDS2615}{Sun} & Environmental exposure and area deprivation vs. paediatric hydrocephalus outcomes & \textit{J Neurosurg Pediatr} & Registry cohort + geospatial & PM$_{2.5}$ (satellite-derived) & USA\\
\addlinespace[0.9mm]
\bottomrule
\end{longtable}}


\vspace{4mm}
\begin{keybox}
{\sffamily\bfseries\footnotesize\color{Deep}Provenance and method}\\[1.2mm]
{\fontsize{7.2}{8.8}\selectfont
Entry window \textbf{4 August 2026}, contiguous with the 3 August issue. Records were
screened on per-record entry date rather than on the harvest window: the sweep ran
4--5 August, but every PubMed record carries an Entrez date of 4 August or earlier and
every Crossref deposit a 4 August \texttt{created} date, so all 13 belong to this issue
and none was held back for the next one.
\textbf{Sources.} PubMed E-utilities and the PubMed connector on two separate axes
(health, instrumentation) windowed on \texttt{[EDAT]} --- 10 health and 2 sensing hits.
Europe PMC \texttt{CREATION\_DATE} --- 2 hits, 1 in scope. Crossref scoped by journal ISSN
across eight aerosol and atmospheric titles that PubMed does not index --- 6 hits, all new
to the corpus, 2 carried. Consensus on four sensing axes (calibration and drift,
satellite--surface fusion, network siting, mobile sensing) --- 3 new records after DOI
resolution and dedup; the free tier caps returns at 3 per sweep, which remains the binding
constraint on the sensing leg. OpenAlex remains blocked behind a paid plan; arXiv
returned nothing in scope after the recency screen.
\textbf{Screening.} Deduplication keys PMID $\rightarrow$ lowercased DOI $\rightarrow$
title signature, checked before summarising. \textbf{Five records were excluded as
off-topic} (agricultural nitrogen accounting, tobacco harm-perception surveys, rice trace
gas fluxes, cough-triggered ventilation, surface-water ecological risk) and \textbf{five
were held, not rejected}: their DOIs verify in Crossref, but no abstract is retrievable
from Crossref, Europe PMC or Consensus and the Elsevier and IEEE landing pages are
client-rendered. Held records are deliberately absent from \texttt{seen.json} so a later
issue can carry them once indexed. Nothing was summarised from a title alone, and no DOI
was inferred from a journal prefix.
\textbf{Labels.} Subtopic, design, particle-metric and geography assignments are
single-label judgements from title, abstract, keywords and MeSH terms, not a controlled
vocabulary. Effect estimates are transcribed verbatim; see the forest-plot caption for
why both are co-pollutant rather than PM estimates.
All DOI links resolve to the publisher of record and were checked by the automated DOI
gate before this issue compiled. Abstracts remain the copyright of their publishers.\par}
\end{keybox}

\end{document}
